%%%%%%%%%%%%%%%%%%%%%%%%%%%%%%%%%%%%%%%%%
% Short Sectioned Assignment LaTeX Template Version 1.0 (5/5/12)
% This template has been downloaded from: http://www.LaTeXTemplates.com
% Original author:  Frits Wenneker (http://www.howtotex.com)
% License: CC BY-NC-SA 3.0 (http://creativecommons.org/licenses/by-nc-sa/3.0/)
%%%%%%%%%%%%%%%%%%%%%%%%%%%%%%%%%%%%%%%%%

% \documentclass[paper=a4, fontsize=11pt]{scrartcl} % A4 paper and 11pt font size
\documentclass[11pt, a4paper]{book}
\usepackage[T1]{fontenc} % Use 8-bit encoding that has 256 glyphs
\usepackage[utf8]{inputenc}
\usepackage{fourier} % Use the Adobe Utopia font for the document - comment this line to return to the LaTeX default
\usepackage{listings} % para insertar código con formato similar al editor
\usepackage[spanish, es-tabla]{babel} % Selecciona el español para palabras introducidas automáticamente, p.ej. "septiembre" en la fecha y especifica que se use la palabra Tabla en vez de Cuadro
\usepackage{url} % ,href} %para incluir URLs e hipervínculos dentro del texto (aunque hay que instalar href)
\usepackage{graphics,graphicx, float} %para incluir imágenes y colocarlas
\usepackage[gen]{eurosym} %para incluir el símbolo del euro
\usepackage{cite} %para incluir citas del archivo <nombre>.bib
\usepackage{enumerate}
\usepackage{hyperref}
\usepackage{graphicx}
\usepackage{tabularx}
\usepackage{booktabs}

\usepackage[table,xcdraw]{xcolor}
\hypersetup{
	colorlinks=true,	% false: boxed links; true: colored links
	linkcolor=black,	% color of internal links
	urlcolor=cyan		% color of external links
}
\renewcommand{\familydefault}{\sfdefault}
\usepackage{fancyhdr} % Custom headers and footers
\pagestyle{fancyplain} % Makes all pages in the document conform to the custom headers and footers
\fancyhead[L]{} % Empty left header
\fancyhead[C]{} % Empty center header
\fancyhead[R]{Lucía González Sánchez} % My name
\fancyfoot[L]{} % Empty left footer
\fancyfoot[C]{} % Empty center footer
\fancyfoot[R]{\thepage} % Page numbering for right footer
%\renewcommand{\headrulewidth}{0pt} % Remove header underlines
\renewcommand{\footrulewidth}{0pt} % Remove footer underlines
\setlength{\headheight}{13.6pt} % Customize the height of the header

\usepackage{titlesec, blindtext, color}
\definecolor{gray75}{gray}{0.75}
\newcommand{\hsp}{\hspace{20pt}}
\titleformat{\chapter}[hang]{\Huge\bfseries}{\thechapter\hsp\textcolor{gray75}{|}\hsp}{0pt}{\Huge\bfseries}
\setcounter{secnumdepth}{4}
\usepackage[Lenny]{fncychap}


\begin{document}

	% Plantilla portada UGR
	\begin{titlepage}
\newlength{\centeroffset}
\setlength{\centeroffset}{-0.5\oddsidemargin}
\addtolength{\centeroffset}{0.5\evensidemargin}
\thispagestyle{empty}

\noindent\hspace*{\centeroffset}\begin{minipage}{\textwidth}

\centering
\includegraphics[width=0.9\textwidth]{logos/logo_ugr.jpg}\\[1.4cm]

\textsc{ \Large TRABAJO FIN DE GRADO\\[0.2cm]}
\textsc{ GRADO EN INGENIERIA INFORMATICA}\\[1cm]

{\Huge\bfseries Heráldica para todas \\}
\noindent\rule[-1ex]{\textwidth}{3pt}\\[3.5ex]
{\large\bfseries Subtítulo }
\end{minipage}

\vspace{2.5cm}
\noindent\hspace*{\centeroffset}
\begin{minipage}{\textwidth}
\centering

\textbf{Autor}\\ {Lucía González Sánchez}\\[2.5ex]
\textbf{Director}\\ {Juan Julián Merelo Guervós}\\[2cm]
\includegraphics[width=0.3\textwidth]{logos/etsiit_logo.png}\\[0.1cm]
\textsc{Escuela Técnica Superior de Ingenierías Informática y de Telecomunicación}\\
\textsc{---}\\
Granada, \today
\end{minipage}
\end{titlepage}


	% Plantilla prefacio UGR
	\thispagestyle{empty}

\begin{center}
{\large\bfseries Heráldica para todas \\ Subtítulo }\\
\end{center}
\begin{center}
Lucía González Sánchez\\
\end{center}

%\vspace{0.7cm}

\vspace{0.5cm}
\noindent\textbf{Palabras clave}: \textit{software libre}
\vspace{0.7cm}

\noindent\textbf{Resumen}\\
	

\cleardoublepage

\begin{center}
	Lucía González Sánchez\\
\end{center}
\vspace{0.5cm}
\noindent\textbf{Keywords}: \textit{open source}, \textit{floss}
\vspace{0.7cm}

\noindent\textbf{Abstract}\\


\cleardoublepage

\thispagestyle{empty}

\noindent\rule[-1ex]{\textwidth}{2pt}\\[4.5ex]

D. \textbf{Tutora/e(s)}, Profesor(a) del ...

\vspace{0.5cm}

\textbf{Informo:}

\vspace{0.5cm}

Que el presente trabajo, titulado \textit{\textbf{Chief}},
ha sido realizado bajo mi supervisión por \textbf{Estudiante}, y autorizo la defensa de dicho trabajo ante el tribunal
que corresponda.

\vspace{0.5cm}

Y para que conste, expiden y firman el presente informe en Granada a Junio de 2018.

\vspace{1cm}

\textbf{El/la director(a)/es: }

\vspace{5cm}

\noindent \textbf{(nombre completo tutor/a/es)}

\chapter*{Agradecimientos}






	% Índice de contenidos
	\newpage
	\tableofcontents

	% Índice de imágenes y tablas
	\newpage
	\listoffigures

	% Si hay suficientes se incluirá dicho índice
	\listoftables 
	\newpage

	% Introducción y descripción del problema
	\chapter{Introducción}

En este capítulo se presenta la motivación del proyecto, así como la definición
del problema que busca abordar. Se incluyen además las historias de usuario que
justifican el desarrollo de la solución propuesta, los objetivos iniciales y la
estructura del documento.

\section{Motivación}
La heráldica es un sistema de representación gráfica codificado que permite
la representación visual de linajes, instituciones y territorios a través de
un conjunto no arbitrario de reglas de diseño y composición. Estas normas
establecidas la hacen un campo apropiado para el análisis estructurado y la
automatización mediante herramientas digitales.

Sin embargo, quienes desean introducirse en esta disciplina se enfrentan a
ciertos obstáculos que dificultan su acceso y comprensión a, debido a que su
documentación se encuentra dispersa y no está completamente digitalizada. Estas 
barreras complican su aprendizaje, aplicación académica y divulgación.

Con la finalidad de satisfacer las necesidades de este grupo de personas y hacer
de esa manera menos hostil el aprendizaje de esta ciencia-arte, este proyecto busca 
desarrollar una herramienta que modernice y agilice el aprendizaje de la heráldica, 
facilitando su comprensión y análisis, aspirando a ofrecer un recurso digital accesible
que contribuya a la preservación y divulgación de esta disciplina.

\section{Definición del problema}

La falta de una herramienta tecnológica que permita analizar y clasificar escudos de 
armas de manera eficiente impide que las personas que estudian o investigan la heráldica
puedan consultar, clasificar o comparar escudos de forma eficiente, ya que en la actualidad
se requieren conocimientos especializados y un acceso a fuentes dispersas, muchas de ellas
no digitalizadas.

Esta situación presenta varias dificultades:

\begin{itemize}
    \item La consulta de escudos requiere revisar múltiples fuentes no centralizadas, 
    en su mayoría analógicas.
    \item No existen bases de datos estandarizadas que permitan buscar escudos por
    características visuales o estructurales.
    \item El análisis se realiza manualmente, lo que ralentiza el proceso y puede
    generar errores.
\end{itemize}

El objetivo de este proyecto será dar solución a los problemas presentados por los
clientes en las siguientes historias de usuario:

\subsection{Historias de usuario}
Con el fin de exponer y comprender mejor las necesidades a las que responde este
proyecto, se presentan los siguientes casos de uso:

\begin{itemize}
    \item \textbf{HU1: Estudiante de Historia del Arte - }Un estudiante de Historia del 
    Arte está paseando por Roma y se encuentra con una obra arquitectónica en cuya fachada
    hay un escudo heráldico. Se pregunta de quién será y qué representa. Sigue caminando 
    y se encuentra con más escudos, lo que le hace pensar en lo oportuna que sería una 
    plataforma que le ayude a identificarlos y obtener información sobre ellos de manera 
    sencilla.
    \item \textbf{HU2: Investigador en heráldica - }Un investigador encuentra tedioso y
    complicado analizar y comprobar escudos en fuentes analógicas, ya que hacerlo manualmente
    de esta manera lo convierte en un proceso lento y propenso a errores.
\end{itemize}

\section{Estructura del proyecto}
Es resto del documento se organiza en los siguientes capítulos:

\begin{itemize}
    \item \textbf{Capítulo 2: Estado del arte - }Se analizan trabajos previos relacionados
    con la digitalización de la heráldica y las herramientas existentes, estableciendo
    el contexto teórico del proyecto.
    \item \textbf{Capítulo 3: Planificación - }Se explica la planificación y metodología
    que se emplean en el desarrollo. Se habla de los milestones del proyecto.
    \item \textbf{Capítulo 4: Análisis - }Define qué debe hacer la herramienta y habla
    sobre su diseño y arquitectura.
    \item \textbf{Capítulo 5: Implementación - }Describe cómo se ha desarrollado y validado
    la herramienta. 
    \item \textbf{Capítulo 6: Conclusiones - }Se resumen los resultados obtenidos y se
    discuten posibles mejoras y líneas futuras de investigación.
\end{itemize}

Este proyecto es software libre, y está liberado con la licencia \cite{gplv3}.

	% Estado del arte
	% 	1. Crítica al estado del arte
	% 	2. Propuesta
	\chapter{Estado del arte}

El software libre y sus licencias \cite{gplv3} ha permitido llevar a cabo una expansión del
aprendizaje de la informática sin precedentes.

	
	\chapter{Planificación}
En esta sección se presenta la planificación del proyecto. Se describen tanto la
metodología empleada como las herramientas elegidas para llevar a cabo el desarrollo,
asegurando la calidad del producto y su entrega continua. Además, se describen las historias
de usuario que exponen las necesidades que este trabajo deberá satisfacer, así como los
primeros hitos que se buscará alcanzar.

\section{Metodología utilizada}
Para la planificación y el desarrollo del proyecto se aplicará un enfoque ágil,
priorizando la entrega continua de valor y la flexibilidad ante los cambios.
En lugar de plantear el desarrollo desde el principio de una manera rígida y detallada,
se definirán hitos, con tareas y objetivos a corto plazo, con entregas frecuentes
y revisiones periódicas para poder detectar y corregir errores a tiempo.

Dentro de este enfoque ágil, se usará una adaptación de SCRUM como marco
de trabajo. Esta metodología ágil estructura el trabajo en ciclos cortos (\textit{sprints}),
de los cuales resultan entregas incrementales del producto. 

Concretamente, cada iteración de este proyecto estará marcada por un hito (\textit{milestone}),
cada uno compuesto por una serie de tareas (\textit{issues}) que se irán completando a lo largo
del desarrollo.

Al final de cada hito, se realizará una comprobación del progreso y se evaluará si cumple con los 
objetivos establecidos, y en caso contrario, se ajustarán los requisitos y se planificarán 
los siguientes hitos.

\subsection{Historias de usuario}
Como punto de partida para la planificación del proyecto, se han definido una seria de 
historias de usuario que reflejan las necesidades de los usuarios que el proyecto final
debe satisfacer. Estas historias de usuario servirán como base para la definición de los hitos.
Las historias de usuario definidas son las siguientes:
\begin{itemize}
    \item \textbf{HU1: Estudiante de Historia del Arte - }Un estudiante de Historia del 
    Arte está paseando por Roma y se encuentra con una obra arquitectónica en cuya fachada
    hay un escudo heráldico papal. Se pregunta de quién será y qué representa. Sigue caminando 
    y se encuentra con más escudos, lo que le hace pensar en lo oportuna que sería una 
    plataforma que le ayude a identificarlos y obtener información sobre ellos de manera 
    sencilla.
    \item \textbf{HU2: Investigador en heráldica - }Un investigador de heráldica española
    encuentra tedioso y complicado analizar y comprobar escudos en fuentes analógicas, ya 
    que hacerlo manualmente de esta manera lo convierte en un proceso lento y propenso a 
    errores.
\end{itemize}

\subsection{Objetivos iniciales}
A partir de las anteriores historias de usuario, se puede establecer el primer
producto mínimamente viable (o primer milestone):
\begin{enumerate}
    \item \textbf{Milestone 0: Infraestructura del proyecto - }Este primer hito tiene como 
    objetivo crear la infraestructura del proyecto, incluyendo la creación del repositorio
    en GitHub, la planificación de los hitos y tareas iniciales, y las herramientas para automatizar
    la compilación y corrección ortográfica y gramatical de la memoria.
    \item \textbf{Milestone 1: Modelo del problema - }Este milestone tiene como objetivo
    elaborar una base sólida mediante la creación de una representación básica y estructurada
    de los escudos, lo que permitirá almacenar y organizar sus atributos clave.
\end{enumerate}

A medida que se avance en este milestone, se establecerá el siguiente paso de manera iterativa,
ajustándose a los resultados obtenidos y con la finalidad de ofrecer valor continuo al usuario.

\subsection{Calidad y entrega continua}
Con el fin de garantizar la calidad del producto y minimizar los errores, se aplicará
un enfoque de control de calidad continuo. Esto implica la implementación de pruebas unitarias
y de integración, así como la revisión y verificación de la documentación generada.
Implementar un enfoque de entrega continua (CD/CI) permitirá automatizar el proceso.


\subsection{Herramientas utilizadas}
Para dar soporte al flujo de trabajo, se utilizarán las siguientes herramientas:
\begin{itemize}
    \item \textbf{Git}: se empleará como sistema de control de versiones para gestionar el código,
    realizar un seguimiento de los cambios, crear ramas de desarrollo y mantener un historial
    claro y ordenado del proyecto.
    \item \textbf{GitHub}: se utilizará como plataforma centralizada para alojar el repositorio
    de código, facilitando la gestión de los \textit{issues} y \textit{milestones}, además de la
    configuración de la integración continua (CI/CD) mediante \textit{GitHub Actions.}
\end{itemize}

\subsubsection{Memoria}
Para la realización de esta memoria, se ha establecido como requisito mínimo que
el documento esté libre de faltas de ortografía y compile correctamente. Partiendo de que 
está escrito en \LaTeX, se han buscado herramientas que permitan la automatización de estas 
comprobaciones. Para ello se han estudiado dos tipos de herramientas: las destinadas a la
corrección gramatical y de estilo, y las empleadas para compilar la memoria.
\\

Entre las \textbf{herramientas para la corrección ortográfica}, han sido analizadas las siguientes
opciones:

\begin{itemize}
    \item \textbf{Textidote}: es una solución de código abierto diseñada específicamente para
    \LaTeX\, lo que evita falsos positivos por malinterpretar los comandos. Requiere Java.
    \item \textbf{LanguageTool}: esta opción de código abierto es más avanzada que la anterior 
    debido a que usa redes neuronales para comprobar la gramática y el estilo. 
    Sin embargo, no está diseñada específicamente para \LaTeX, lo que hace necesario 
    invertir tiempo extra en eliminar los comandos antes de realizar las verificaciones.
    Tiene una versión gratuita y una versión prémium con funcionalidades avanzadas.
    \item \textbf{GNU Aspell}: es el corrector ortográfico estándar de GNU, aunque compila para
    otros sistemas operativos. Es de código abierto, no corrige gramática ni estilo, solo
    ortografía, y no está diseñado específicamente para \LaTeX.
\end{itemize}

Después de la comparativa, la opción elegida ha sido \textbf{Textidote}, porque a pesar de 
no ser tan avanzada en la corrección, supone un ahorro considerable de tiempo el hecho de no 
tener que eliminar los comandos de \LaTeX\ al verificar la gramática y el estilo, como sucede
con las otras dos alternativas analizadas.
\\

Con respecto a las \textbf{herramientas para la compilación de la memoria}, se han estudiado:

\begin{itemize}
    \item \textbf{TeX Live}: es la distribución por defecto de \LaTeX\ para Linux. Sin embargo,
    la instalación puede ser pesada y excesiva si únicamente se necesita compilar documentos
    sencillos.
    \item \textbf{Overleaf}: es un editor \LaTeX\ basado en la nube, por lo que no requiere 
    instalación previa. Es ideal para trabajar de manera colaborativa al permitir
    a varios usuarios trabajar en el mismo documento y ver los cambios a tiempo real.
    No obstante, tiene limitaciones en cuanto a capacidad, que pueden solventarse con
    la versión de pago.
    \item \textbf{TinyTex}: es una versión basada en Tex Live, pero mucho más ligera. Está orientada
    sobre todo a usuarios de R, aunque es posible instalar manualmente los paquetes 
    específicos que sean necesarios, lo que la convierte en una buena opción en un contexto 
    en el que los recursos están limitados.
\end{itemize}

Dado a la complejidad añadida de instalar manualmente los paquetes específicos necesarios
para el proyecto y que los recursos no son limitantes en este caso, la herramienta elegida
es \textbf{TeX Live}. Se descarta \textbf{Overleaf} debido a la limitación de la versión gratuita,
además de que no será necesario un entorno de trabajo colaborativo.

\section{Temporización}
La temporización se establecerá en función de los hitos y las tareas necesarias para 
alcanzar los objetivos propuestos.

	% Análisis del problema
	% 1. Análisis de requisitos
	% 2. Análisis de las soluciones
	% 3. Solucion propuesta
	% 4. Análisis de seguridad
	\chapter{Análisis del problema}
Este capítulo tiene como objetivo analizar en profundidad el problema al que se pretende 
dar solución, abordando tanto los aspectos conceptuales del dominio como las implicaciones 
técnicas y funcionales del sistema desde una perspectiva inicial. No se busca definir una
solución, sino comprender el contexto del problema y acotarlo de manera realista.

Se analizarán los riesgos técnicos asociados y se esbozarán posibles vías de evolución futura. 
Finalmente, se elegirán el lenguaje de programación y las herramientas que servirán de base 
para las decisiones de implementación que se describirán en el capítulo siguiente.

\section{Análisis de riesgos}
Antes de comenzar el desarrollo se identifican los siguientes riesgos:

\begin{itemize}
    \item \textbf{Ambigüedad conceptual: } ciertos términos heráldicos pueden tener distintos
    significados según la tradición.
    \item \textbf{Falta de datos estructurados: } no existen fuentes públicas completas con escudos
    de forma estandarizada y digital, lo que puede dificultar la extracción automática de información.
\end{itemize}

\section{Evolución futura}
Aunque no forma parte del alcance inicial y siguiendo la metodología el desarrollo es incremental, 
es importante considerar desde un principio ciertos aspectos que podrían evolucionar en el futuro:

\begin{itemize}
    \item \textbf{Integración de inteligencia artificial: } para identificar visualmente los escudos.
    Puede ser un objetivo a largo plazo relacionado con la HU1.
    \item \textbf{Ampliación del sistema: } a otras tradiciones heráldicas o idiomas.
\end{itemize}

Contemplar estas posibilidades desde el inicio permite que la selección de herramientas y tecnologías
no limite o complique el crecimiento futuro del sistema.

\section{Selección de herramientas}
En esta sección se detallan las herramientas y tecnologías que se han considerado para el desarrollo
del software, así como los criterios que se han seguido para su selección. La selección de estas
herramientas se ha basado en las secciones previas.

\subsection{Lenguaje de programación}\label{sec:lenguaje-programación}
Antes de avanzar en el diseño del modelo de datos y la arquitectura del sistema, es necesario definir 
el lenguaje de programación, ya que cada lenguaje ofrece estructuras, herramientas y recursos diferentes 
para representar y manipular la información.

Para seleccionar el lenguaje más adecuado, se han identificado una serie de requisitos técnicos derivados
de las historias de usuario:

\begin{itemize}
    \item Debe permitir el desarrollo de una plataforma accesible desde distintos dispositivos 
    (móviles y ordenadores) y, preferiblemente, con capacidad multiplataforma o web. (\hyperref[sec:hu1]{HU1})
    \item Debe facilitar búsquedas rápidas de información. (\hyperref[sec:hu1]{HU1}, \hyperref[sec:hu2]{HU2})
    \item Se valora la posibilidad de funcionamiento offline para situaciones sin conexión. (\hyperref[sec:hu1]{HU1})
    \item Es recomendable que permita integrar técnicas de procesamiento de imagen y funcionalidades de inteligencia artificial. (\hyperref[sec:hu1]{HU1})
    \item Debe facilitar la gestión de grandes volúmenes de datos y ofrecer mecanismos para su análisis, filtrado y clasificación. (\hyperref[sec:hu2]{HU2})
\end{itemize}

A partir de estos requisitos, se han definido los siguientes criterios de búsqueda y selección.


\subsubsection{Criterios de búsqueda}
Para la selección de lenguajes de programación se han considerado los siguientes criterios:

\begin{itemize}
    \item \textbf{Curva de aprendizaje moderada: } que me permita como desarrolladora aprender
    rápidamente y ser productiva desde el principio.
    \item \textbf{Comunidad activa:} que asegure soporte continuo, documentación actualizada y 
    disponibilidad de librerías externas.
    \item \textbf{Adecuación al dominio del problema:} especialmente en lo relativo a 
    manipulación de datos estructurados y representación semántica.
    \item \textbf{Flexibilidad:} posibilidad de evolucionar el proyecto hacia otros ámbitos, 
    como el desarrollo web o la integración de inteligencia artificial.
\end{itemize}

\subsubsection{Criterios de selección}
Para la selección final del lenguaje de programación se han considerado los siguientes criterios:

\begin{itemize}
    \item Facilidad de aprendizaje y sintaxis clara.
    \item Ecosistema y disponibilidad de librerías.
    \item Comunidad y soporte continuos.
    \item Facilidad de despliegue y mantenimiento.
\end{itemize}

\subsubsection{Lenguajes considerados}
Entre los lenguajes de programación que se han considerado para el desarrollo del software, se 
incluyen los siguientes:

\begin{itemize}
    \item \textbf{JavaScript: }
    \begin{itemize}
        \item \textbf{Curva de aprendizaje baja:} sintaxis sencilla
        y recursos de aprendizaje abundantes.
        \item \textbf{Ecosistema:} muchas librerías y frameworks
        para desarrollo web.
        \item \textbf{Comunidad activa:} comunidad enorme y
        soluciones para casi todo.
        \item \textbf{Soporte continuo:} actualizaciones y mejoras
        constantes.
        \item \textbf{Facilidad de despliegue:} se ejecuta en
        navegador o servidor con Node.js.
    \end{itemize}

    \item \textbf{TypeScript: }extensión de JavaScript tipada.
    \begin{itemize}
        \item \textbf{Curva de aprendizaje:} algo mayor que JS,
        pero aporta robustez.
        \item \textbf{Ecosistema:} mismo ecosistema que JS,
        compatible con librerías.
        \item \textbf{Comunidad activa:} gran comunidad y
        uso creciente en proyectos grandes.
        \item \textbf{Soporte continuo:} soporte de Microsoft,
        mejoras constantes.
        \item \textbf{Facilidad de despliegue:} se compila a
        JavaScript y se despliega igual.
    \end{itemize}
        
    \item \textbf{Python: }
        \begin{itemize}
        \item \textbf{Curva de aprendizaje baja:} sintaxis clara y sencilla, 
        mucha documentación y ejemplos.
        \item \textbf{Ecosistema:} enorme cantidad de librerías para 
        manejo de datos y automatización.
        \item \textbf{Comunidad activa:} comunidad muy grande y foros 
        con soluciones para problemas comunes.
        \item \textbf{Soporte continuo:} tiene soporte estable y 
        actualizaciones frecuentes.
        \item \textbf{Facilidad de despliegue:} se despliega bien 
        en servidores y servicios cloud.
    \end{itemize}

    \item \textbf{Rust: }
    \begin{itemize}
        \item \textbf{Curva de aprendizaje:} más alta que Python
        o JS por su enfoque bajo nivel.
        \item \textbf{Ecosistema:} en crecimiento, aunque menos
        librerías para datos.
        \item \textbf{Comunidad activa:} comunidad activa y
        muy colaborativa.
        \item \textbf{Soporte continuo:} soporte sólido y
        versiones estables.
        \item \textbf{Facilidad de despliegue:} binarios ligeros
        y fáciles de ejecutar.
    \end{itemize}

    \item \textbf{Go: }
    \begin{itemize}
        \item \textbf{Curva de aprendizaje:} sintaxis simple,
        fácil de aprender.
        \item \textbf{Ecosistema:} buenas librerías para web
        y servidores.
        \item \textbf{Comunidad activa:} comunidad sólida,
        buena documentación.
        \item \textbf{Soporte continuo:} soporte fuerte de Google,
        actualizaciones regulares.
        \item \textbf{Facilidad de despliegue:} genera binarios,
        muy fácil de desplegar.
    \end{itemize}

    \item \textbf{C++: }
    \begin{itemize}
        \item \textbf{Curva de aprendizaje:} alta, sintaxis más
        compleja y detallada. Sin embargo, es el lenguaje que 
        mejor conozco.
        \item \textbf{Ecosistema:} librerías maduras, aunque menos
        cómodas para web y datos.
        \item \textbf{Comunidad activa:} comunidad grande,
        documentación extensa.
        \item \textbf{Soporte continuo:} soporte estable,
        estándar bien definido.
        \item \textbf{Facilidad de despliegue:} requiere más
        configuración y compilación.
    \end{itemize}

\end{itemize}

\subsubsection{Lenguaje seleccionado}
De la anterior lista, compararé ahora solo TypeScript, Python y Go, ya que son los que
mejor se adaptan a los criterios de selección y a mis necesidades como desarrolladora.

\begin{itemize}
    \item \textbf{Curva de aprendizaje:} Python es el más fácil por su sintaxis clara y
    natural, siendo famoso por ser lenguaje de iniciación a la programación.
    \item \textbf{Ecosistema y librerías:} TypeScript tiene un ecosistema muy amplio 
    para desarrollo web, mientras que Python destaca en ciencia de datos y automatización.
    Go es más limitado.
    \item \textbf{Comunidad activa:} TypeScript y Python tienen comunidades enormes y
    activas, Go es más nicho.
    \item \textbf{Soporte continuo:} los tres tienen soporte activo, aunque Python gana en
    cuanto a trayectoria.
    \item \textbf{Facilidad de despliegue:} Go es muy fácil de desplegar, ya que produce
    binarios que no dependen de entornos externos. TypeScript requiere Node.js
    y Python necesita un intérprete, lo que puede complicar el despliegue en algunos
    entornos.
\end{itemize}

En conclusión, después de comparar directamente los lenguajes entre ellos, es evidente
que tanto TypeScript como Python son opciones muy sólidas y podría elegir cualquiera de las
dos. Sin embargo, atendiendo a posibles líneas de evolución futura, he decidido optar por
\textbf{Python}, ya que ofrece mayor versatilidad para trabajar con datos estructurados y 
permite, a medio y largo plazo, integrar funcionalidades de inteligencia artificial sin 
grandes cambios tecnológicos.

\subsection{Gestor de dependencias y del entorno}\label{sec:gestor-dependencias}
Para poder garantizar que el proyecto sea idempotente, mantenible, idempotente y fácil de instalar en 
cualquier entorno, es fundamental contar con un gestor de dependencias y del entorno que
siga los estándares oficiales de Python. 

Seguir los estándares oficiales de Python significa que la herramienta aplica las normas y recomendaciones
definidas en los PEP (\textit{Python Enhancement Proposals}), que son documentos oficiales que describen 
nuevas funcionalidades, mejoras o guías de estilo en Python. En particular, los PEP relevantes para la
gestión de dependencias son:

\begin{itemize}
    \item \textbf{\href{https://peps.python.org/pep-0508/}{PEP 508: }} Especifica la sintaxis 
    para declarar dependencias en Python (tales como nombre del paquete, versión y extras).
    \item \textbf{\href{https://peps.python.org/pep-0518/}{PEP 518: }} Define cómo los paquetes 
    Python deben declarar sus dependencias necesarias para ejecutar el sistema de construcción. y en este 
    caso también se aplica a proyectos de tipo aplicación como el nuestro para garantizar la idempotencia 
    y reproducibilidad del entorno. Además, introduce el archivo \texttt{pyproject.toml} como lugar 
    centralizado para estas especificaciones.
\end{itemize}

Esto permitirá instalar y documentar todas las dependencias necesarias de forma automática,
evitando la instalación manual y reduciendo el riesgo de errores o incompatibilidades. 

\subsubsection{Criterios de búsqueda}
Se necesita un gestor de paquetes que:
\begin{itemize}
    \item Soporte los estándares actuales de Python (PEP 508 y 518).
    \item Sea adecuada para un proyecto pequeño, priorizando simplicidad sobre complejidad 
    innecesaria.
\end{itemize}
Estas condiciones limitan las opciones a herramientas modernas que implementen \texttt{pyproject.toml},
excluyendo herramientas usadas habitualmente como \texttt{pip} + \texttt{venv} o \texttt{pipenv}.

\subsubsection{Criterios de selección}
Entre las opciones que cumplen los requisitos anteriores, se prioriza la velocidad; la herramienta 
debe resolver dependencias y crear o actualizar entornos rápidamente, minimizando el tiempo de 
instalación y configuración.

\subsubsection{Herramientas consideradas}
Se han comparado tres opciones:
\begin{itemize}
    \item \textbf{Poetry}: Implementa nativamente los PEP 508 y 518, gestionando dependencias
    y entorno con \texttt{pyproject.toml}. Asegura que las instalaciones sean reproducibles y
    ofrece una resolución de dependencias rápida y fiable.
    \item \textbf{Flit:} Se centra en la creación de paquetes y su distribución, utiliza
    \texttt{pyproject.toml} para definir dependencias y es muy rápido. Sin embargo, no gestiona
    entornos virtuales, lo que supone varios problemas: las dependencias se instalan  en el sistema
    global, no quedan aisladas, lo que puede causar conflictos entre proyectos; rompe la idempotencia,
    ya que las instalaciones pueden variar según el entorno del sistema; y complica la automatización
    de tests y despliegue, puesto que no asegura un entorno consistente. Por lo tanto, esta herramienta por
    sí sola no es suficiente para proyectos que requieren un entorno controlado.
    \item \textbf{Hatch:} También soporta \texttt{pyproject.toml} y ofrece gestión de entornos
    y dependencias. Es flexible y potente, aunque puede ser lenta en proyectos grandes. 
    \item \textbf{uv: } Herramienta escrita en \texttt{Rust}, mucho más rápida que las demás.
    Sin embargo, su adopción es aún limitada y su ecosistema es mucho más reducido que el de
    \texttt{Poetry}, lo que podría dificultar su integración.
    \item \textbf{PDM: } Se denomina como una mezcla entre \texttt{pip}, \texttt{poetry} y 
    \texttt{virtualenv}, sin embargo, no es tan popular ni tiene tanta comunidad como \texttt{poetry}.
    Está basado en \texttt{\href{https://peps.python.org/pep-0582/}{PEP 582}}, que propone una alternativa a los entornos virtuales tradicionales,
    utilizando el directorio \texttt{\_\_pypackages\_\_} en lugar de \texttt{venv}. Sin embargo,
    este PEP no está tan extendido ni es tan compatible con otras herramientas del ecosistema
    Python como los PEP 508 y 518.
\end{itemize}

\subsubsection{Herramienta seleccionada}
Se ha seleccionado \textbf{Poetry} por ajustarse mejor a los estándares actuales de Python, ya
que emplea \texttt{pyproject.toml}, definido en el \textit{PEP 518}, de manera nativa para la
gestión de dependencias y soporta la especificación de versiones y extras según el \textit{PEP 508}.

Además, Poetry garantiza la idempotencia y reproducibilidad del entorno, ofreciendo el mejor equilibrio
entre estabilidad, madurez, velocidad y soporte de la comunidad.

\subsection{Herramienta de testeo} \label{sec:herramienta-testeo}
Para garantizar la calidad del código y poder aplicar desarrollo guiado por pruebas 
(\textit{Test Driven Development, TDD}), es necesario que se seleccione una herramienta de testeo antes de 
comenzar a modelar el dominio. Esta práctica, propuesta inicialmente por Kent Beck \cite{beckTDD}, 
consiste en escribir primero las pruebas y después el código que las satisface, lo que permite asegurar 
que el sistema es fiable desde sus fases iniciales. De esta forma, cada entidad y objeto de valor podrá 
validarse desde el inicio mediante pruebas automatizadas.

Python ya incluye en su biblioteca estándar el módulo \texttt{unittest}, considerado la herramienta
de testeo por defecto. Sin embargo, su sintaxis es algo verbosa y obliga a escribir mucho código
\textit{boilerplate} (es decir, código repetitivo) para definir pruebas sencillas, lo que puede 
retrasar el desarrollo.

Es por esta razón por la que se opta por usar \texttt{pytest}, una herramienta muy popular que actúa 
como \textit{test runner}, es decir, se encarga de ejecutar las pruebas. \texttt{pytest} no sustituye 
a \texttt{\href{https://docs.python.org/es/3/library/unittest.html}{unittest}}, sino que lo utiliza por 
debajo, lo que permite que se mantenga la compatibilidad total con él. Esto significa que cualquier 
prueba escrita con \texttt{unittest} funcionará perfectamente con \texttt{pytest}, siendo ambas 
complementarias y no excluyentes.

\texttt{pytest} cuenta con un amplio ecosistema de plugins que amplían sus funcionalidades, y aunque no son
estrictamente necesarios, hay algunos que podrían resultar útiles, como \texttt{pytest-cov}, que hace
de puente entre \texttt{pytest} y \texttt{\href{https://coverage.readthedocs.io/en/7.10.6/}{coverage.py}}
(herramienta estándar), y permite medir la cobertura de código durante los test.

\subsection{Verificación automática de buenas prácticas}
El código no solamente debe funcionar, sino ser \textbf{legible}, \textbf{consistente} y \textbf{mantenible},
según los principios del llamado \emph{Clean Code}~\cite{martin2009cleancode}.
Aplicar buenas prácticas desde el inicio evita errores triviales y agiliza las revisiones. Para reforzarlo,
este proyecto incorpora comprobación automática que se ejecuta en cada push o pull request (CI).

\subsubsection{Análisis estático y formateo automático}
Para mantener la calidad del código, se emplea análisis estático (\emph{linting}) que detecta tempranamente 
errores y malas prácticas, y formateo automático que garantiza un estilo consistente conforme a 
\href{https://peps.python.org/pep-0008/}{PEP 8}, mejorando legibilidad y mantenibilidad.

Las herramientas consideradas han sido:

\begin{itemize}
  \item \textbf{Ruff}: muy rápido; reimplementa reglas habituales de \texttt{pyflakes}/\texttt{pycodestyle}
  (errores comunes y estilo), ordena imports y aplica modernizaciones sencillas. Integra el formateador, por
  lo que requiere poca configuración.
  \item \textbf{Flake8}: linter clásico y estable. Flexible, pero suele requerir combinar
  varios plugins, lo que aumenta configuración y tiempo de ejecución.
  \item \textbf{Pylint}: linter con análisis más exhaustivo (convenciones y algunos olores de diseño). Tiende
  a ser más verboso y lento; puede generar más ruido en proyectos pequeños.
  \item \textbf{Black}: es el formatter más usado por la comunidad; resultados muy estables y predecibles. Implica
  utilizar una herramienta adicional además del linter.
  \item \textbf{autopep8 / yapf}: linters históricos hoy menos extendidas.
\end{itemize}

\subsubsection{Herramienta seleccionada}
Se selecciona \textbf{Ruff} como herramienta única para lint y formateo, debido a su ejecución rápida y su menor 
complejidad operativa (es una sola herramienta).
Se integra en CI con un workflow que se lanza en cada push o pull request y hace fallar la integración si hay 
errores de estilo o de formato.

\subsection{Base de datos}\label{sec:base-datos}
Antes de comenzar con el diseño de la base de datos, es necesario definir las herramientas que se emplearán para 
su creación y gestión. En este caso, no basta con elegir una sola herramienta, sino que deben seleccionarse dos 
componentes distintos aunque complementarios: por un lado, el motor de base de datos, que es el responsable de 
almacenar y gestionar los datos y, por otro lado, la capa de acceso a datos, que es la interfaz que permite 
interactuar con la base de datos desde el código Python, traduciendo las clases y objetos del dominio a tablas y
registros en la base de datos.

Esta separación de las herramientas permite mantener el proyecto modular y flexible, facilitando cambios futuros sin
que ello afecte a toda la arquitectura del sistema.

\subsubsection{Motor de base de datos}
El motor de base de datos es el componente encargado de almacenar físicamente los datos, gestionar las consultas y
mantener la integridad de la información. En este proyecto se ha considerado necesario utilizar un motor relacional,
ya que permite representar de manera natural las relaciones entre las entidades del dominio.

En este proyecto no se han considerado bases de datos NoSQL o documentales, debido a que el dominio no presenta las 
características que justifican su uso. Los datos heráldicos tienen una estructura fija y fuertemente relacional: 
un escudo está compuesto por un campo, uno o varios muebles y, opcionalmente, un adorno exterior, todos ellos con 
relaciones y restricciones bien definidas. Este tipo de información se modela de manera más natural y segura en un 
sistema relacional, donde las dependencias y reglas de integridad pueden expresarse explícitamente.

\paragraph{Herramientas consideradas}
Se han evaluado las siguientes opciones:

\begin{itemize}
    \item \textbf{SQLite:} Es una base de datos ligera que guarda toda la información en un único archivo local. 
    No requiere servidor ni configuración adicional, lo que facilita su utilización en entornos de desarrollo y pruebas.
    Soporta el estándar SQL completo y ofrece un rendimiento más que suficiente para el tamaño de datos esperado
    durante la fase de desarrollo. Además, es totalmente compatible con los principales ORM de Python.

    \item \textbf{PostgreSQL:} Es un sistema de base de datos relacional orientado a entornos de producción.
    Ofrece gran robustez, soporte para consultas avanzadas y manejo eficiente de múltiples usuarios concurrentes.
    Sin embargo, requiere la instalación y configuración de un servidor, lo que introduce una complejidad innecesaria
    para un entorno de desarrollo local.

    \item \textbf{DuckDB:} Motor de base de datos analítico moderno, muy eficiente para consultas sobre grandes volúmenes
    de datos, pero menos adecuado para operaciones transaccionales o persistencia tradicional de registros, que es
    el caso de este proyecto.
\end{itemize}

\paragraph{Herramienta seleccionada}
Para este proyecto se selecciona \textbf{SQLite} como motor de base de datos.  
Su simplicidad, portabilidad y compatibilidad con el ecosistema Python la hacen ideal para el desarrollo incremental
guiado por pruebas que sigue este trabajo. Al almacenar los datos en un único archivo, evita dependencias externas
y permite crear y destruir la base de datos en cada iteración sin afectar al resto del sistema, lo que facilita las
pruebas unitarias y el desarrollo local.

Además, en caso de que el proyecto evolucione hacia una aplicación web o de mayor escala, la transición a un motor
más robusto como PostgreSQL puede realizarse sin grandes cambios en el código gracias al uso de una capa de acceso
independiente.

\subsubsection{Capa de acceso a datos}
La capa de acceso a datos es la responsable de conectar el código Python con la base de datos, permitiendo trabajar 
con las tablas mediante clases y objetos en lugar de escribir consultas SQL manuales. Este enfoque se conoce como 
\textit{Object-Relational Mapping} (ORM) \cite{fowlerPatterns}, y facilita que el modelo de datos del programa 
refleje directamente la estructura del dominio, manteniendo la coherencia con las clases ya definidas en el proyecto.

\paragraph{Herramientas consideradas}
Se han valorado las siguientes opciones:

\begin{itemize}
    \item \textbf{SQLAlchemy ORM:} Es la más usada en el ecosistema Python. Permite un control muy 
    detallado sobre las consultas y las relaciones entre tablas, además de ofrecer una gran flexibilidad para 
    personalizar el comportamiento del modelo. Sin embargo, su sintaxis es algo más extensa y requiere escribir 
    más código para operaciones simples, lo que puede ralentizar el desarrollo de un prototipo.

    \item \textbf{SQLModel:} Es una librería moderna basada en SQLAlchemy, pero con una sintaxis más concisa y 
    tipada, inspirada en las \texttt{dataclasses} de Python. Permite definir modelos con anotaciones de tipo, 
    haciendo que el código sea más legible y coherente con el estilo ya empleado en este proyecto. Mantiene 
    compatibilidad total con SQLAlchemy y puede evolucionar sin cambios si se requiere un entorno más complejo.

    \item \textbf{Peewee:} Es un ORM ligero y sencillo, apropiado para proyectos pequeños. Ofrece una sintaxis 
    más directa, no obstante cuenta con una comunidad más reducida y menor integración con herramientas modernas de 
    validación y tipado, lo que limitaría su escalabilidad a medio plazo.

    \item \textbf{Django ORM:} Incluye un potente sistema de gestión de base de datos, sin embargo, está diseñado para 
    funcionar dentro del framework web Django. Adoptarlo implicaría introducir dependencias y una estructura 
    de proyecto completamente diferente, sin aportar ventajas reales en este contexto.
\end{itemize}

\paragraph{Herramienta seleccionada}
Para este proyecto se ha elegido \textbf{SQLModel} como capa de acceso a datos.  
Ofrece una sintaxis sencilla y moderna, basada en anotaciones de tipo, que encaja bien con las clases del dominio 
ya desarrolladas en los hitos anteriores. Esto permite definir las tablas de la base de datos de forma muy similar 
a las \texttt{dataclasses} que ya se utilizan, reduciendo la necesidad de aprender una sintaxis completamente nueva.

Además, SQLModel está construido sobre \texttt{SQLAlchemy}, por lo que mantiene su solidez y compatibilidad con 
motores como SQLite o PostgreSQL. Esta combinación facilita el desarrollo y deja abierta la posibilidad de ampliar 
el sistema en el futuro sin tener que reescribir la lógica actual.

\subsection{Herramienta de scraping}\label{sec:scraping}
Para extraer datos de páginas estáticas (Wikipedia) se han considerado tres opciones:

\begin{itemize}
  \item \textbf{Requests + BeautifulSoup:} Esta combinación usa \texttt{Requests} para realizar las peticiones 
  HTTP/S a la fuente, obteniendo el HTML, y \texttt{BeautifulSoup} para parsear el HTML y extraer la información 
  relevante. Es una solución ligera, rápida de implementar y adecuada para páginas estáticas como Wikipedia y
  proyectos pequeños.
  \item \textbf{Scrapy:} framework robusto de alto rendimiento con gestión de colas, middlewares y pipelines.
  Aporta más complejidad y estructura de proyecto de la que se necesita para este trabajo.
  \item \textbf{Playwright/Selenium:} orientados a páginas dinámicas con JavaScript; sobrecarga innecesaria para 
  Wikipedia, que sirve HTML estático.
\end{itemize}

\subsubsection{Herramienta seleccionada}
Se ha seleccionado \textbf{Requests + BeautifulSoup} por su simplicidad, velocidad de desarrollo y ajuste al alcance 
del proyecto: permite iterar rápido, añade poca dependencia y cubre las necesidades del scraping planteado.

\subsection{Herramientas de backend}
El backend actúa como intermediario entre la base de datos de blasones y la interfaz web. 
Su función es recibir las peticiones del usuario, procesarlas y devolver la respuesta correspondiente.  
Para ello se requiere una herramienta que facilite la gestión de rutas y peticiones sin tener que implementar 
manualmente todo el protocolo HTTP.

Entre las opciones consideradas se encuentran:

\begin{itemize}
    \item \textbf{Django} es un framework de tipo \textit{MVC} (Modelo-Vista-Controlador). 
    Proporciona una estructura completa que incluye ORM, autenticación y panel de administración, 
    pensada para desarrollar aplicaciones web de gran tamaño. En este proyecto resultaría excesivo, 
    ya que solo se necesita un servidor que gestione unas pocas rutas y renderice plantillas sencillas.
    \item \textbf{FastAPI} es un framework orientado al desarrollo de servicios \textit{REST}, 
    diseñado para construir APIs que devuelven datos en formato JSON. Su rendimiento y tipado lo hacen 
    ideal para aplicaciones basadas en microservicios, pero no está centrado en generar páginas HTML ni 
    en servir contenido visual.
    \item \textbf{Flask} es un \textit{microframework} minimalista, pensado para crear aplicaciones web 
    pequeñas de forma flexible. Permite definir rutas, procesar peticiones y renderizar plantillas, sin 
    imponer una estructura rígida ni dependencias adicionales.
\end{itemize}

Aunque estas tres herramientas permiten desarrollar aplicaciones web en Python, no son directamente 
comparables entre sí. Cada una responde a un enfoque de desarrollo diferente: \textbf{Django} sigue el 
patrón \textit{MVC} \cite{reenskaug1979mvc} e incluye todos los componentes de una aplicación completa; \textbf{FastAPI} se 
orienta a arquitecturas \textit{REST} \cite{fielding2000rest}, donde el backend actúa como proveedor de datos para otros servicios; 
y \textbf{Flask} se sitúa en un punto intermedio, como un microframework que permite combinar ambos 
enfoques según las necesidades del proyecto.  

En este caso, el sistema no requiere un modelo complejo ni un servicio API, sino una capa ligera que 
conecte los datos con la interfaz web y renderice las vistas. Por ello, \textbf{Flask} resulta la opción 
más adecuada, ya que ofrece el equilibrio necesario entre simplicidad, flexibilidad y funcionalidad.

\subsection{Herramientas para la interfaz de usuario} \label{sec:interfaz-herramientas}
Con el objetivo de garantizar una implementación sencilla, reproducible y coherente con la arquitectura modular del sistema, 
se analizaron distintas herramientas y librerías para el desarrollo de la interfaz web.



\subsubsection{Gestión de peticiones y controladores}
Cuando el usuario realiza una búsqueda, la aplicación debe recibir esa solicitud, procesarla y devolver los 
resultados adecuados. Para ello conviene usar un framework que gestione las rutas y las peticiones web,
proporcionando una estructura base sin tener que implementarlo todo desde cero. Esto agiliza el desarrollo y
reduce la probabilidad de errores.

Se evaluaron tres opciones de frameworks: \textit{Django}, \textit{FastAPI} y \textit{Flask}. 
\textit{Django} fue descartado porque impone una estructura demasiado grande para un proyecto de este tamaño, 
y \textit{FastAPI} porque está más pensado para crear servicios API que para renderizar páginas, por lo que
se eligió \textbf{Flask}, ya que es ligero, muy fácil de integrar con el modelo de datos del sistema y 
suficiente para el tipo de interfaz que se busca desarrollar.

\subsubsection{Generación de vistas}
Una vez obtenidos los resultados, es necesario mostrarlos en una página HTML. 
Para no generar el código HTML manualmente, se utiliza un motor de plantillas, 
que permite generar una estructura base y rellenarla con los datos de cada escudo.  
Flask incluye de forma nativa \textbf{Jinja2}, un motor sencillo y potente que evita añadir más dependencias y
permite construir páginas reutilizables, mantener el código limpio y separar claramente la presentación visual 
de la lógica del programa.

\subsubsection{Actualización dinámica de contenido}
Aunque la interfaz es muy simple, resulta más cómodo que el usuario pueda ver los resultados 
sin tener que recargar toda la página cada vez que realiza una búsqueda.  
En lugar de usar un framework de JavaScript completo (como React o Vue), que añadiría mucha complejidad 
innecesaria, se eligió \textbf{HTMX}, una pequeña biblioteca que permite actualizar solo una parte de la página 
(por ejemplo, la lista de resultados) mediante peticiones asíncronas al servidor. De este modo se mejora la 
experiencia de uso sin abandonar la estructura ligera del proyecto.

\subsubsection{Presentación visual}
Por último, es importante que la interfaz resultara clara y agradable sin invertir tiempo 
en diseñar una hoja de estilos desde cero.  
Para ello se utiliza \textbf{Pico.css}, una hoja de estilos que se carga directamente desde Internet y aplica un diseño 
limpio y accesible. De esta forma no es necesario instalar ni configurar librerías adicionales.

\subsubsection{Herramientas seleccionadas}
El conjunto de herramientas escogido ofrece un equilibrio entre simplicidad técnica y una buena experiencia de usuario. 
Permite tener una aplicación funcional, fácilmente reproducible y adaptable a futuras mejoras, como la inclusión de 
filtros adicionales o módulos de inteligencia artificial para clasificar los escudos.

	% Desarrollo bajo sprints: 
	% 	1. Permitir registros y login de usuarios
	% 	2. Desarrollo del sistema de incidencias
	% 	3. Desarrollo del sistema de denuncias administrativas y accidentes
	% 	4. Desarrollo del sistema de croquis
	%   5. Instalación de la aplicación de manera automática
	\chapter{Implementación}

La implementación del software se ha dividido en hitos. Estos han sido definidos en GitHub
y cada uno de ellos contiene un grupo de \textit{issues} que se corresponden con las distintas
mejoras que se han ido incorporando al software a lo largo de su desarrollo.

Cada milestone representa un bloque de trabajo con un objetivo definido, lo que permite seguir
un enfoque ágil, enfocado en construir progresivamente un prototipo sólido en lugar de intentar
desarrollar todo el sistema de una vez. 

\section{Milestone 1: Elección del lenguaje de programación}
En este primer hito se decidió el lenguaje de programación a utilizar para el desarrollo del software,
siendo Python la opción elegida. Esta elección se justificó y documentó debidamente en el
\hyperref[sec:lenguaje-programación]{anterior capítulo}.

\section{Milestone 2: Modelo del problema}
Este hito se centra en dos tareas principales: la configuración de las herramientas necesarias para
el desarrollo guiado por pruebas, y la definición del modelo del dominio siguiendo el enfoque 
\textit{Domain-Driven Design} (DDD).

\subsection{Configuración del entorno de desarrollo}
En el capítulo anterior se seleccionaron las herramientas que se emplearían para 
\hyperref[sec:gestor-dependencias]{gestionar las dependencias y el entorno} y 
\hyperref[sec:herramienta-testeo]{realizar pruebas}. En este hito se procedió a su instalación y 
configuración.

\subsubsection{Poetry}
Se ha usado \textit{Poetry} para gestionar las dependencias y el entorno virtual del proyecto.

Al inicializar el gestor, se genera automáticamente un archivo \texttt{pyproject.toml} en la raíz del
repositorio, que contiene la información básica del proyecto y servirá como punto central
para la gestión de dependencias.

\subsubsection{Pytest}
La herramienta seleccionada para las pruebas unitarias ha sido \texttt{pytest}.
La elección de esta herramienta ya fue justificada en el capítulo anterior.

Se ha añadido \texttt{pytest} como dependencia, lo que significa que se instalará únicamente
para el entorno de desarrollo y no se considera una dependencia necesaria para la ejecución del 
prototipo en producción. De esta manera se mantiene una separación clara entre las librerías 
empleadas para programar y probar el software, y aquellas que serán estrictamente necesarias para 
su despliegue.

Al instalar dependencias mediante Poetry, junto al archivo \texttt{pyproject.toml} se genera 
automáticamente el archivo \texttt{poetry.lock}. El primero declara qué dependencias se requieren y 
en qué rango de versiones son aceptables; en cambio, el archivo \texttt{poetry.lock} registra la 
versión exacta de esas dependencias en el momento de la instalación. Gracias a este mecanismo se 
garantiza la reproducibilidad del entorno, ya que cualquier persona que clone el repositorio e 
instale las dependencias obtendrá exactamente las mismas versiones usadas en el desarrollo 
original.

\subsection{Modelo del dominio}
El modelado del dominio se ha realizado siguiendo el enfoque \textit{Domain-Driven Design} (DDD)
propuesto por Eric Evans \cite{evansDDD}. Este método de diseño de software se centra en que
el sistema refleje fielmente la realidad del dominio que se está modelando, evitando que los 
aspectos técnicos o de implementación condicionen la representación del problema.

Uno de los principios fundamentales de DDD es la creación de un \textit{lenguaje ubicuo}, introducido
por Evans en el capítulo 2 de su obra \cite{evansDDD}. Esto consiste en definir términos y conceptos 
específicos del dominio para que los conceptos empleados en el código coincidan con los del propio 
dominio. De este modo se facilita la comunicación entre el análisis y la implementación. En este 
proyecto se ha definido el siguiente lenguaje ubicuo:

\begin{itemize}
    \item \textbf{Escudo}: Entendemos como escudo a la representación gráfica de un conjunto de
    elementos heráldicos que identifican a una persona, entidad o institución.
    \item \textbf{Blasón}: Descripción textual del escudo.
    \item \textbf{Portador}: Persona, entidad o institución a la que pertenece el escudo.
    \item \textbf{Campo}: Superficie del escudo donde se disponen los elementos heráldicos.
    \item \textbf{Boca}: Forma exterior del escudo, es decir, su contorno.
    \item \textbf{Partición}: Cada una de las divisiones del campo del escudo. 
    \item \textbf{Esmalte}: Color o metal con el que se pinta un elemento heráldico.
    \item \textbf{Pieza heráldica}: Son elementos de distinto esmalte al del campo que suelen tocar
    los bordes del escudo.
    \item \textbf{Figura o mueble}: Son todo aquello que se encuentra dentro del campo del escudo y
    no resulta ser una pieza heráldica.
    \item \textbf{Adorno exterior}: Son todos aquellos elementos que no se encuentran dentro del campo
    del escudo.
\end{itemize}

\begin{figure}[H]
    \centering
    \includegraphics[width=1\textwidth]{figuras/escudo-papa-francisco.png}
    \caption{Partes del escudo de armas del papa Francisco \cite{wikiEscudo}.}
    \label{fig:partes-escudo}
\end{figure}

Las definiciones empleadas en este lenguaje ubicuo se han extraído del 
\textit{Manual de heráldica: la ciencia del blasón} \cite{delgadoHeraldica}.

En esta primera iteración se ha decidido limitar el modelo a una parte reducida del lenguaje ubicuo. 
Concretamente, se trabajará con el \textbf{Escudo} como composición de otros elementos, incorporando 
inicialmente el \textbf{Campo} y el \textbf{Esmalte}, que constituyen la base mínima necesaria para 
definir un prototipo funcional. 

Otros elementos como la \textbf{Boca}, el \textbf{Portador}, las \textbf{Particiones}, las 
\textbf{Piezas heráldicas}, las \textbf{Figuras} o los \textbf{Adornos exteriores} se han recogido 
en el lenguaje ubicuo, pero quedan fuera del alcance de esta iteración y se abordarán en futuras 
extensiones del modelo.

\subsection{Desarrollo dirigido por pruebas (TDD)}

El modelo inicial del dominio se ha implementado siguiendo la metodología de 
\textit{Test-Driven Development} (TDD) \cite{beckTDD}. Esto significa que antes de escribir 
el código se definieron las pruebas que debía superar, y después se implementó lo mínimo necesario 
para que pasaran.

En esta primera iteración se definieron pruebas unitarias para los elementos básicos 
del modelo: \texttt{Esmalte}, \texttt{Campo} y \texttt{Escudo}. Dichas pruebas validan:

\begin{itemize}
    \item Que un \texttt{Escudo} válido puede crearse a partir de un \texttt{Campo} 
    con un \texttt{Esmalte} permitido.
    \item Que los intentos de producir un \texttt{Esmalte} con valores no reconocidos producen 
    una excepción.
    \item Que no es posible instanciar un \texttt{Campo} sin un \texttt{Esmalte} válido.
\end{itemize}

Al principio estas pruebas no pasan, ya que no existe ninguna implementación. A continuación se
escribe el código mínimo necesario para que todas las pruebas definidas pasen correctamente. Esto
se repetirá cada vez que se añada una nueva funcionalidad al modelo.

Como resultado surgieron las clases \texttt{Esmalte}, \texttt{Campo} y \texttt{Escudo}, que constituyen 
la base mínima del dominio:

\begin{itemize}
    \item \texttt{Esmalte}: representa los colores y metales permitidos en heráldica. 
    Comprueba que el valor introducido sea válido y lo normaliza.
    \item \texttt{Campo}: representa la superficie del escudo. Solo puede generarse 
    si se le proporciona un \texttt{Esmalte} válido.
    \item \texttt{Escudo}: actúa como elemento principal y se compone de un \texttt{Campo}.
\end{itemize}

\section{Milestone 3: Identificación y filtrado por esmalte}

En este milestone se introduce la primera lógica de negocio, que consiste en filtrar los escudos
por el esmalte del campo. Esto facilita la identificación de escudos de \hyperref[sec:hu1]{HU1} y 
la clasificación y el filtrado de escudos de \hyperref[sec:hu2]{HU2}.

Para ello, es necesario implementar un catálogo mínimo de ejemplo para poder realizar la consulta.

Los criterios de aceptación son sencillos: (1) dado un esmalte válido, la consulta devuelve 
únicamente escudos con ese campo; (2) sin criterio, se devuelve el catálogo completo; y 
(3) con un criterio erróneo, la respuesta es una lista vacía.

\subsection{Adición de un catálogo mínimo}

En este punto no había ningún escudo cargado y, por tanto, no había nada que filtrar. Antes 
de implementar la búsqueda por esmalte del campo, lo razonable era añadir al sistema un catálogo 
base sobre el que poder trabajar y comprobar el comportamiento.

Para este hito opté por un catálogo mínimo y en memoria, dentro del propio paquete. 
Cada entrada se modela como una pequeña ficha con el nombre del escudo y el esmalte.
El campo se guarda como un Esmalte, de modo que queda validado y normalizado desde
el minuto uno y se evitan incoherencias.

Se han añadido a este catálogo siete escudos, uno para cada esmalte canónico, con la finalidad de
poder probar la búsqueda de cualquier esmalte. La decisión de mantener el catálogo en memoria
simplifica la implementación inicial y permite centrarse en la funcionalidad de filtrado sin
introducir complejidades adicionales que se tratarán en futuros hitos, como la implementación
de una base de datos.

Siguiendo TDD, primero escribí unas pruebas muy básicas: que el catálogo no esté vacío y que el 
campo de cada ficha sea un Esmalte. Con esas pruebas en rojo implementé lo mínimo para ponerlas en
verde: un catálogo pequeño en memoria y la función listar(). Con eso ya está lista la base para el
siguiente paso: filtrar la búsqueda por el esmalte del campo.

	% Presupuesto

	% Conclusiones
	\chapter{Conclusiones y trabajos futuros}



	% Trabajos futuros


	
	\newpage
	\bibliography{bibliografia}
	\bibliographystyle{plain}
	
\end{document}

